\section{Event-clock studies, monitors, and tactical results}
\label{sec:tactical}

\paragraph{The clock itself.}
Entering each name on the day its information became public (per
holder-filing acceptance) instead of the D+45 batch adds
$+0.21\%$/yr on the long side ($t=0.18$) and $+0.31\%$/yr on the
short ($t=-0.22$): the cost of the batch clock is measured at
approximately zero, consistent with the delay distribution of
Section~\ref{sec:delays}.

\paragraph{Classical interactions: five pre-registered conditional
effects, none confirmed.}
Theory motivates conditioning each 13F signal on a classical
characteristic; five such interactions were pre-registered with
target cells, expected signs, and a Bonferroni bar of $|t|\ge2.57$
for the family. None confirmed (Table~\ref{tab:inter}); NCB
replicated unconditionally at $t=+3.78$ inside the same module ---
the factor is broad, not a corner-case effect.

\begin{table}[t]
\centering\footnotesize
\caption{Pre-registered 13F$\times$classical interactions, 50
quarters. Verdict statistic: paired delta between the target-tercile
spread and the unconditional spread, NW(4).}
\label{tab:inter}
\begin{tabular}{llrrl}
\toprule
interaction (theory) & target & uncond.\ (t) & delta (t) & verdict \\
\midrule
NCB $\times$ mktcap (limits to arbitrage) & small & $+3.78$ & $+1.34$ & not confirmed \\
dbreadth $\times$ IO (\citealt{nagel2005}-style) & low IO & $+1.17$ & $-1.42$ & not confirmed \\
FSS $\times$ ADV (impact where illiquid) & illiquid & $-0.90$ & $+1.05$ & not confirmed \\
pressure $\times$ vol (flow moves price in high vol) & high vol & $+1.09$ & $+0.77$ & not confirmed \\
exit\_rate $\times$ momentum (exodus in losers) & losers & $+0.96$ & $+2.02$ & wrong sign \\
\bottomrule
\end{tabular}
\end{table}

\paragraph{The pressure autopsy: why flow pressure should pay, why
it mostly does not, and where the lesson survived.}
The economic case for a flow-pressure factor is genuine, and its
failure is as informative as NCB's success. A 13F reveals flows that
were \emph{already executed}: by the decision date the quarter's
selling has largely happened and its price impact is realized inside
the filing quarter --- which is why the unconditional pressure
signal dies once orthogonalized to momentum, reversal, and the
concurrent quarter's return; there is nothing left to forecast. The
one situation where pressure should still pay \emph{forward} is an
\emph{unfinished liquidation}: a position large relative to ADV
cannot be unwound in one quarter, so the filing reveals pending
supply rather than consumed supply. This is measured, not
conjectured: the mechanism stage of Section~\ref{sec:fss} shows
distressed managers selling 34.5\% of remaining shares in the
quarter \emph{after} the filing --- liquidation demonstrably
continues past the disclosure.

Two conditioning routes follow, with different fates. Route one
conditions on the \emph{stock}: pressure should matter more where
volatility is high (valuation uncertainty amplifies the impact of
uninformed flow) --- the E4 interaction above, and the project's
cleanest example of a flicker dying properly: promising in smoke,
$t=+2.06$ on the old universe (just under the promotion bar), an
identification autopsy --- conditioning on residual volatility
versus ADV versus market cap to separate valuation uncertainty from
disguised illiquidity --- that returned ``exploratory, no formal
claim'' ($t=-0.50/-0.61$), and a fall to $t=+0.77$ on complete
books: non-promotion vindicated by data it had never seen. Route two
conditions on the \emph{seller}: restrict to managers whose selling
is \emph{forced}, where continuation is guaranteed by the redemption
itself rather than inferred from the stock's characteristics. That
route is FSS, and it is the one that prices. The design lesson:
\textbf{condition on the agent's constraint, not on the asset's
state} --- the same asymmetry as Law 1.

\paragraph{Imbalance reversal (\citealt{miori2022}).}
Replicated in oracle form and re-evaluated with the biases removed
(point-in-time snapshots instead of perfect foresight of the
cross-section; universe-churn control): the adjusted premium is a
fraction of the published point. We report both numbers as the
clearest example in the project of how much 13F results depend on
who-knew-what-when accounting.

\paragraph{Quarter-end effects.}
The window-dressing design (learned manager propensity, within-loser
test to kill the momentum confound, permutation placebo) rejected
both pre-registered signs, with the placebo tracking the thesis ---
whatever exists lives in ownership/ADV structure, not in
window-dressing propensity. A residual, honestly labeled
hypothesis-generating: the post-photo reversal is entirely a
December phenomenon ($-1.91\%$ in Q4 versus $-0.05\%$ otherwise,
eleven annual observations), observationally equivalent to
portfolio pumping \citep{carhart2002} or a January
low-institutional-ownership bounce; with $n=11$ it remains a
footnote.

\paragraph{Amendment events.}
Across 250k direction-signed events, material restatements move
prices \emph{against} the direction of the correction
($-0.10\%$ at $+3$d, $t=-3.52$; clean pre-event placebo);
administrative amendments are zero, as they must be; confidential
NEW-HOLDINGS reveals drift positively at $+10$d but with a
contaminated placebo. The pre-registered thesis (price follows the
correction) is dead; the inverted fact is consistent with amendments
following stale-price attention rather than creating news.

\paragraph{Daily monitors.}
Three daily rules were evaluated against pre-registered bars on the
drifted-book machinery. The short-squeeze alarm inverted into the
useful rule: shorts spiking $+7\%$/5d on doubled volume
\emph{revert} ($-1.21\%$ over the next 5d versus $-0.12\%$ calm,
$t=-2.60$) --- the anti-panic rule is to hold, not cover.
The aggregate deleveraging index is a risk conditioner, not a
signal: correlation $+0.14$ with future ten-day market volatility,
zero with future returns (Law 1 again); throttling gross above its
expanding 80th percentile improves the combo proxy's Sharpe
0.49$\to$0.55 and drawdown $-17\%\to-15\%$. Early covering at
apparent exhaustion loses $1.5\%$/quarter to save 7 bps of borrow
--- the pressure continues; the short is held through the quarter.
The post-exhaustion tactical \emph{long} died at its gate
(stress doubles the forced-sale rate, $12.2\%$ vs $6.7\%$, but the
thesis-minus-placebo event return is $t=+0.95$, and the only
significant cell is \emph{continuation}): distressed books do not
V-bounce.
